\begin{frame}{Transport Layer: TCP over Multi-Hop Wireless}

\begin{center}
\textbf{TCP's sending rate is approximately $cwnd/RTT$ and is paced by returning acknowledgements.}
\end{center}

\vspace{0.02cm}

\begin{columns}[T,onlytextwidth]
\column{0.48\textwidth}
\fontsize{9}{10.8}\selectfont
\textbf{TCP control loop}

\begin{center}
\begin{tikzpicture}[
    node distance=0.20cm,
    stage/.style={
        draw=ResearchColor,
        rounded corners=3pt,
        fill=CardBg,
        minimum width=5.25cm,
        minimum height=0.58cm,
        align=center,
        font=\fontsize{9}{10.5}\selectfont
    },
    flow/.style={-{Latex[length=2.1mm]},draw=ResearchColor,line width=0.75pt}
]
    \node[stage] (cwnd) {$cwnd$ limits unacknowledged data};
    \node[stage,below=of cwnd] (inject) {Packets traverse every wireless hop};
    \node[stage,below=of inject] (feedback) {ACK clock controls new transmissions};

    \draw[flow] (cwnd) -- (inject);
    \draw[flow] (inject) -- (feedback);
\end{tikzpicture}
\end{center}

\column{0.48\textwidth}
\fontsize{9}{10.8}\selectfont
\textbf{Why multi-hop wireless is difficult}
\begin{itemize}
    \item Every hop consumes shared channel time.
    \item Neighboring hops contend; MAC retries vary the RTT.
    \item Wireless loss may look like congestion to TCP.
    \item A larger $cwnd$ intensifies MAC contention.
\end{itemize}
\end{columns}

\end{frame}
